\section{Geographic Determinants}\label{sec:geo}

This section employs linguistic cartography and a focused literature review to
examine how geography has shaped the distribution and divergence of Philippine
languages, combining the visualisation of geolinguistic variation with the
historical and sociocultural context that accounts for it.

Each of the three regions introduced in \cref{sec:intro} possesses distinct
ecological and historical conditions corresponding to linguistic
variation~\cite{Boquet_2017}, a diversity scholars attribute largely to the
archipelago's fragmented geography, which has historically restricted mobility
and sustained linguistic isolation~\cite{Blust_1991}. Mountains, rivers, and
ocean channels act as both connectors and separators. The \emph{barangay}, the
smallest socio-political unit of pre-colonial society, was typically sited near
rivers or coasts, locations that facilitated trade while simultaneously creating
micro-environments of linguistic autonomy in which neighbouring islands could
remain mutually unintelligible despite their proximity.

Blust's Austronesian dispersal hypothesis~\cite{Blust_1991} holds that a single
protolanguage, introduced by early settlers from the Asian mainland, diverged
through localised adaptation and restricted contact, accumulating phonological
and lexical innovations as communities became isolated and yielding the 175
documented Philippine languages, two of which are now
extinct~\cite{ethnologue2025}. The sixteen studied here, located
in~\cref{sec:intro}, range from dominant regional lingua francas to endangered
local tongues.

\subsection{Methodology}

We employed two official linguistic atlases published by the Komisyon ng Wikang
Filipino (KWF), the \emph{Mapa ng mga Wika} and the \emph{Mapa ng mga Wika ng
Katutubong Pamayanang Kultural}, which delineate the spatial distribution of
languages across Philippine administrative boundaries and identify overlapping
zones of multilingual interaction~\cite{Marrion2025MapaKultural,
    KWF2025MapaPilipinas}. To enhance precision, we generated a custom digital map
from boundary data published by the United Nations Office for the Coordination of
Humanitarian Affairs (OCHA), rendered with Matplotlib and GeoPandas in
Python~\cite{OCHA_Philippines_2018_PHL_Admin_Boundaries, Hunter:2007,
    kelsey_jordahl_2020_3946761}, ensuring spatial consistency between the
linguistic and geographic datasets.

Topographic data, comprising elevation models, hydrology, and terrain
classifications, were obtained from the National Mapping and Resource
Information Authority (NAMRIA)~\cite{NAMRIA_Downloads}. The analysis proceeded by
geospatial alignment of the linguistic and physical layers, detection of
convergence between linguistic boundaries and terrain discontinuities, and
qualitative interpretation through historical and sociolinguistic literature,
with particular attention to how elevation and water systems delineate the
boundaries of the languages in \cref{tab:ph_languages}.

\subsection{Results and Discussion}

The generated linguistic map, presented in \cref{fig:lingmap_out}, reveals clear
patterns of correspondence between physical terrain and language boundaries.
Ilocano and Paranan, for instance, are divided by the Cordillera and Sierra Madre
ranges, while the seas separating Panay, Negros, and Leyte correspond precisely to
the divisions among Hiligaynon, Kinaray-a, and Waray. The case of Bikol,
Masbatenyo, and Romblomanon illustrates an intermediate linguistic zone where
mutual intelligibility correlates with ease of maritime travel.

In contrast, lowland regions such as Central Luzon and Northern Mindanao show
diffuse linguistic borders, reflecting trade, migration, and urbanisation. The
presence of Chavacano in Zamboanga exemplifies contact-induced change, in which
Spanish influence merged with Cebuano and Tausug substrates to form a stable
creole. Yami, situated beyond the archipelago's northern limit, correspondingly
preserves continuity with the languages of Taiwan through geographic proximity
and historical seafaring ties.

The spatial correlations therefore underscore geography as both a limiting and an
enabling force: rugged terrain and isolated islands encourage linguistic
diversification, while navigable waterways foster contact and hybrid formation.
These findings align with global typologies of geography-driven diversity,
placing the Philippines alongside regions such as Indonesia and Papua New Guinea
in the broader Austronesian context.

